\chapter{极限与连续}\label{ch:g12:limcont}

本章将\omterm{def:g12:seq:sequence}{数列}极限的概念推广到
实变量\omterm{def:g11:func:function}{函数}，并引入\omterm{def:g12:limcont:continuity}{连续}性。核心结果是
中间值定理，它把定性信息（“$f$ \omterm{def:g12:limcont:continuity}{连续}
且变号”）转化为\omterm{def:g10:algebra:equation}{方程}解的存在性。

\section{函数的极限}

下文中，$f$ 是定义在\omterm{def:g10:numbers:interval}{区间}或若干\omterm{def:g10:numbers:interval}{区间}的\omterm{def:g10:numbers:interunion}{并集}
$D \subseteq \R$ 上的\omterm{def:g11:func:function}{函数}。

\begin{definition}[无穷远处的极限]\label{def:g12:limcont:liminf}
设 $f$ 定义在\omterm{def:g10:numbers:interval}{区间} $\intco{a}{+\infty}$ 上，$\ell \in \R$。
称 $f$ \emph{趋于 $\ell$}\index{极限!函数}（在 $+\infty$ 处），
记作 $\lim\limits_{x\to+\infty} f(x) = \ell$，若
包含 $\ell$ 的每个开\omterm{def:g10:numbers:interval}{区间}都包含 $f(x)$ 在 $x$ 充分大时的全部取值：
即对任意 $\varepsilon > 0$，存在 $A$，使得当 $x \geq A$ 时
有 $\abs{f(x) - \ell} \leq \varepsilon$。

称 $f$ 在 $+\infty$ 处趋于 $+\infty$，若对任意 $M \in \R$，
存在 $A$，使得 $f(x) \geq M$ 对所有 $x \geq A$ 成立。
在 $-\infty$ 处的极限及极限为 $-\infty$ 的情形可类似定义。
\end{definition}

\begin{definition}[一点处的极限]\label{def:g12:limcont:limpoint}
设 $a \in \R$，$f$ 在 $a$ 附近有定义（$a$ 处可除外）。
称 $f$ 趋于 $\ell \in \R$（在 $a$ 处），记作
$\lim\limits_{x \to a} f(x) = \ell$，若对任意 $\varepsilon > 0$，
存在 $\delta > 0$，使得对所有满足 $x \in D$ 且
$\abs{x - a} \leq \delta$ 时，都有 $\abs{f(x) - \ell} \leq \varepsilon$。

称 $f$ 趋于 $+\infty$（在 $a$ 处），若对任意 $M$，存在 $\delta > 0$，
使得 $f(x) \geq M$ 在 $x \in D$ 且 $\abs{x - a} \leq \delta$ 时成立。
单侧极限（$x \to a^+$、$x \to a^-$）把 $x$ 限制为 $x > a$
或 $x < a$。
\end{definition}

\begin{example}
$\lim\limits_{x\to+\infty} \frac{1}{x} = 0$,
$\lim\limits_{x\to 0^+} \frac{1}{x} = +\infty$,
$\lim\limits_{x\to 0^-} \frac{1}{x} = -\infty$。\omterm{def:g11:func:function}{函数}
$x \mapsto \frac1x$ 在 $0$ 处没有（双侧）极限。
\end{example}

\begin{definition}[渐近线]\label{def:g12:limcont:asymptote}
直线 $y = \ell$ 是\emph{水平渐近线}\index{渐近线}，
即 $f$ \omterm{def:g10:functions:graph}{图像}在 $+\infty$（或 $-\infty$）处，若
$\lim_{x\to+\infty} f(x) = \ell$（或在 $-\infty$ 处）。直线 $x = a$ 是
\emph{铅直渐近线}\index{渐近线}，若 $f$ 在 $a$ 处有无穷的单侧极限。
\end{definition}

\begin{omfigure}
\begin{tikzpicture}
\begin{axis}[omaxis,
  width=8.8cm, height=6cm,
  xmin=-2.9, xmax=4.9, ymin=-5.2, ymax=9.2,
  xtick={1}, ytick={2}]
  \draw[omProp, dashed] (axis cs:-2.9,2) -- (axis cs:4.9,2);
  \draw[omProp, dashed] (axis cs:1,-5.2) -- (axis cs:1,9.2);
  \addplot[omDef, thick, domain=-2.8:0.855, samples=60] {2 + 1/(x-1)};
  \addplot[omDef, thick, domain=1.145:4.8, samples=60] {2 + 1/(x-1)};
\end{axis}
\end{tikzpicture}

{\small $f(x) = 2 + \frac{1}{x-1}$ 的\omterm{def:g10:functions:graph}{图像}，及其\omterm{def:g12:limcont:asymptote}{水平
渐近线} $y = 2$（在 $\pm\infty$ 处）与\omterm{def:g12:limcont:asymptote}{铅直渐近线} $x = 1$。}
\end{omfigure}

\begin{proposition}[极限的运算]\label{prop:g12:limcont:operations}
\cref{prop:g12:seq:operations} 中的运算法则（和、积、商）
对\omterm{def:g11:func:function}{函数}的极限（一点处或无穷远处）同样成立，
并有相同的四种未定式 $(+\infty)+(-\infty)$、$0\times\infty$、
$\frac{\infty}{\infty}$、$\frac{0}{0}$。
\end{proposition}

\begin{proof}
证明与\omterm{def:g12:seq:sequence}{数列}情形完全相同，只需把“对 $n \geq N$”
换成“对 $x \geq A$”或“对 $\abs{x-a} \leq \delta$”。
\end{proof}

\begin{theorem}[复合函数的极限]\label{thm:g12:limcont:composition}
设 $f, g$ 为\omterm{def:g11:func:function}{函数}，$a, b, c$ 为\omterm{def:g10:numbers:sets}{实数}或 $\pm\infty$。若
$\lim\limits_{x \to a} f(x) = b$，且 $\lim\limits_{y \to b} g(y) = c$，则
\[
\lim_{x \to a} g\bigl(f(x)\bigr) = c .
\]
把 $f$ 换成\omterm{def:g12:seq:sequence}{数列}时，结论同样成立：若
$u_n \to b$，且 $\lim_{y\to b} g(y) = c$，则 $g(u_n) \to c$。
\end{theorem}

\begin{proof}
不妨设 $b, c$ 为有限数（其他情形类似）。任给 $\varepsilon > 0$。
存在 $\delta > 0$，使得由 $\abs{y - b} \leq \delta$ 可推出
$\abs{g(y) - c} \leq \varepsilon$；又存在 $\delta' > 0$，使得
由 $\abs{x - a} \leq \delta'$ 可推出 $\abs{f(x) - b} \leq \delta$。
两式相接，即得 $\abs{g(f(x)) - c} \leq \varepsilon$，
对 $\abs{x - a} \leq \delta'$ 成立。
\end{proof}

\begin{theorem}[比较与夹逼]\label{thm:g12:limcont:squeeze}
设 $f, g, h$ 为\omterm{def:g11:func:function}{函数}，$a \in \R \cup \{\pm\infty\}$。
\begin{enumerate}
  \item 若 $f \leq g$ 在 $a$ 附近，且 $f(x) \to +\infty$（$x \to a$），
        则 $g(x) \to +\infty$。
  \item 若 $f \leq g \leq h$ 在 $a$ 附近，且 $f, h$ 趋于同一个有限极限
        $\ell$（在 $a$ 处），则 $g(x) \to \ell$（$x \to a$）。
\end{enumerate}
\end{theorem}

\begin{proof}
与\omterm{def:g12:seq:sequence}{数列}情形的论证相同（见\cref{thm:g12:seq:squeeze}）。
\end{proof}

\begin{method}[极限的计算方法]\label{met:g12:limcont:limits}
\begin{itemize}
  \item 多项式与有理\omterm{def:g11:func:function}{函数}在 $\pm\infty$ 处的极限：提出
        $x$ 的最高次幂，
        极限即为首项之比的极限。
  \item $\frac{0}{0}$ 型：在点 $a$ 处，从分子与分母中提出因式 $(x-a)$，
        或识别出一个\omterm{def:g11:deriv:derivative}{导数}
        $\lim_{x\to a}\frac{f(x)-f(a)}{x-a} = f'(a)$（见\cref{ch:g12:deriv}）。
  \item 含平方根时：乘以共轭表达式。
  \item 振荡因子（$\cos$、$\sin$）：使用夹逼定理。
\end{itemize}
\end{method}

\section{连续性}

\begin{definition}[连续性]\label{def:g12:limcont:continuity}
设 $f$ 为定义在\omterm{def:g10:numbers:interval}{区间} $I$ 上的\omterm{def:g11:func:function}{函数}，$a \in I$。称 $f$ 在
$a$ 处\emph{连续}\index{连续}，是指 $\lim\limits_{x \to a} f(x) = f(a)$。
称该\omterm{def:g11:func:function}{函数}在 $I$ 上\emph{连续}\index{连续}，
若该\omterm{def:g11:func:function}{函数}在 $I$ 的每一点处都连续。
\end{definition}

直观上，\omterm{def:g10:numbers:interval}{区间}上\omterm{def:g12:limcont:continuity}{连续}\omterm{def:g11:func:function}{函数}的\omterm{def:g10:functions:graph}{图像}
可以不抬笔一笔画出。

\begin{proposition}[常见函数的连续性]\label{prop:g12:limcont:usual}
多项式、有理\omterm{def:g11:func:function}{函数}、$x \mapsto \sqrt{x}$、$x \mapsto \abs{x}$、
$\cos$、$\sin$、$\exp$ 与 $\ln$ 在其\omterm{def:g11:func:function}{定义域}内所含的任一\omterm{def:g10:numbers:interval}{区间}上\omterm{def:g12:limcont:continuity}{连续}。
\omterm{def:g12:limcont:continuity}{连续}\omterm{def:g11:func:function}{函数}的和、积、商（在有定义处）及复合
仍为\omterm{def:g12:limcont:continuity}{连续}\omterm{def:g11:func:function}{函数}。
\end{proposition}

\begin{proof}[部分证明]
和、积、商及复合的稳定性由极限的相应定理得出。
恒等映射 $x \mapsto x$ 与常值\omterm{def:g11:func:function}{函数}的\omterm{def:g12:limcont:continuity}{连续}性
可直接由定义得到，因此多项式与有理\omterm{def:g11:func:function}{函数}\omterm{def:g12:limcont:continuity}{连续}。
其余常见\omterm{def:g11:func:function}{函数}的\omterm{def:g12:limcont:continuity}{连续}性在它们被构造的章节中证明，
或在本阶段不加证明地承认。
\end{proof}

\begin{example}
\emph{取整\omterm{def:g11:func:function}{函数}} $x \mapsto \floor{x}$
（即满足 $\leq x$ 的最大\omterm{def:g10:numbers:sets}{整数}）在任一\omterm{def:g10:numbers:sets}{整数} $n$ 处都不\omterm{def:g12:limcont:continuity}{连续}：
其在该点的左极限为 $n - 1$，而 $\floor{n} = n$。
\end{example}

\begin{omfigure}
\begin{tikzpicture}
\begin{axis}[omaxis,
  width=7.2cm, height=5.6cm,
  xmin=-1.9, xmax=3.5, ymin=-2.5, ymax=2.7,
  xtick={-1,1,2,3}, ytick={-2,-1,1,2}]
  \draw[omDef, thick] (axis cs:-1.5,-2) -- (axis cs:-1,-2);
  \draw[omDef, thick] (axis cs:-1,-1) -- (axis cs:0,-1);
  \draw[omDef, thick] (axis cs:0,0) -- (axis cs:1,0);
  \draw[omDef, thick] (axis cs:1,1) -- (axis cs:2,1);
  \draw[omDef, thick] (axis cs:2,2) -- (axis cs:3,2);
  \fill[omDef] (axis cs:-1,-1) circle (1.5pt);
  \fill[omDef] (axis cs:0,0) circle (1.5pt);
  \fill[omDef] (axis cs:1,1) circle (1.5pt);
  \fill[omDef] (axis cs:2,2) circle (1.5pt);
  \draw[omDef, fill=white] (axis cs:-1,-2) circle (1.5pt);
  \draw[omDef, fill=white] (axis cs:0,-1) circle (1.5pt);
  \draw[omDef, fill=white] (axis cs:1,0) circle (1.5pt);
  \draw[omDef, fill=white] (axis cs:2,1) circle (1.5pt);
  \draw[omDef, fill=white] (axis cs:3,2) circle (1.5pt);
\end{axis}
\end{tikzpicture}

{\small 取整\omterm{def:g11:func:function}{函数}在每个\omterm{def:g10:numbers:sets}{整数}处都有跳跃：\omterm{def:g11:func:function}{函数}值（实心点）
与左极限（空心点）不同。}
\end{omfigure}

\begin{proposition}[数列与连续性]\label{prop:g12:limcont:seqcont}
若 $u_n \to \ell$ 且 $f$ 在 $\ell$ 处\omterm{def:g12:limcont:continuity}{连续}，则 $f(u_n) \to f(\ell)$。
特别地，若由 $u_{n+1} = f(u_n)$ 所定义的\omterm{def:g12:seq:sequence}{数列}\omterm{def:g12:seq:limit}{收敛}于
$\ell$，且 $f$ 在 $\ell$ 处\omterm{def:g12:limcont:continuity}{连续}，则 $f(\ell) = \ell$。
\end{proposition}

\begin{proof}
第一个结论由对\omterm{def:g12:seq:sequence}{数列} $(u_n)$ 应用\cref{thm:g12:limcont:composition}得出。
对于第二个结论，在 $u_{n+1} = f(u_n)$ 两边取极限：
左边趋于 $\ell$，右边趋于 $f(\ell)$，
而极限是唯一的。
\end{proof}

\section{介值定理}

\begin{theorem}[介值定理]\label{thm:g12:limcont:ivt}
\index{介值定理}
设 $f$ 在 $\intcc{a}{b}$ 上\omterm{def:g12:limcont:continuity}{连续}，$k$ 为介于
$f(a)$ 与 $f(b)$ 之间的任一数。则存在 $c \in \intcc{a}{b}$，使得
$f(c) = k$。
\end{theorem}

\begin{omfigure}
\begin{tikzpicture}
\begin{axis}[omaxis,
  width=8.8cm, height=6cm,
  xmin=-0.4, xmax=4.4, ymin=-0.5, ymax=4.7,
  xtick={0.5,3.19,3.8}, xticklabels={$a$,$c$,$b$},
  ytick={1.26,2.5,3.98}, yticklabels={$f(a)$,$k$,$f(b)$}]
  \addplot[omProp, dashed, domain=0:3.19] {2.5};
  \addplot[omDef, thick, domain=0.5:3.8] {0.08*x^3 - 0.5*x + 1.5};
  \draw[black, dashed, thin] (axis cs:0.5,0) -- (axis cs:0.5,1.26)
    -- (axis cs:0,1.26);
  \draw[black, dashed, thin] (axis cs:3.8,0) -- (axis cs:3.8,3.98)
    -- (axis cs:0,3.98);
  \draw[black, dashed, thin] (axis cs:3.19,2.5) -- (axis cs:3.19,0);
  \fill (axis cs:0.5,1.26) circle (1.5pt);
  \fill (axis cs:3.8,3.98) circle (1.5pt);
  \fill (axis cs:3.19,2.5) circle (1.5pt);
\end{axis}
\end{tikzpicture}

{\small 介值定理：一条\omterm{def:g12:limcont:continuity}{连续}曲线连接
$(a, f(a))$ 与 $(b, f(b))$，必与每条水平直线
$y = k$ 相交，这些直线的高度介于 $f(a)$ 与 $f(b)$ 之间。}
\end{omfigure}

\begin{proof}
将 $f$ 替换为 $f - k$（必要时替换为 $k - f$），不妨设
$f(a) \leq 0 \leq f(b)$，并寻找 $f$ 的一个零点。
用\emph{二分法}\index{二分法}构造两个\omterm{def:g12:seq:sequence}{数列}：取 $a_0 = a$，$b_0 = b$；
已知 $a_n \leq b_n$ 且 $f(a_n) \leq 0 \leq f(b_n)$，令
$m = \frac{a_n + b_n}{2}$，并定义
\[
(a_{n+1}, b_{n+1}) =
\begin{cases}
(a_n, m) & \text{ 若 } f(m) \geq 0,\\
(m, b_n) & \text{ 若 } f(m) < 0.
\end{cases}
\]
两种情形下都有 $f(a_{n+1}) \leq 0 \leq f(b_{n+1})$，\omterm{def:g12:seq:sequence}{数列} $(a_n)$
\omterm{def:g12:seq:monotonic}{递增}，$(b_n)$ \omterm{def:g10:functions:variations}{递减}，且
$b_n - a_n = \frac{b-a}{2^n} \to 0$：两\omterm{def:g12:seq:sequence}{数列}相邻（见
\cref{exo:g12:seq:10}），并\omterm{def:g12:seq:limit}{收敛}于公共极限
$c \in \intcc{a}{b}$。由\omterm{def:g12:limcont:continuity}{连续}性，$f(a_n) \to f(c)$ 且
$f(b_n) \to f(c)$；在 $f(a_n) \leq 0$ 与
$f(b_n) \geq 0$ 中取极限，得 $f(c) \leq 0 \leq f(c)$，从而 $f(c) = 0$。
\end{proof}

\begin{remark}
\omterm{thm:g12:limcont:ivt}{二分法}证明是一种算法：它可给出任意精确的
解的\omterm{def:g10:numbers:approx}{近似值}，每一步误差减半。这正是计算机
便能求解 $f(x) = 0$，只需会计算 $f$ 的值。
\end{remark}

\begin{theorem}[双射定理]\label{thm:g12:limcont:bijection}
设 $f$ 在\omterm{def:g10:numbers:interval}{区间} $I$ 上\omterm{def:g12:limcont:continuity}{连续}且\emph{严格\omterm{def:g12:seq:monotonic}{单调}}。则
对于每个 $k$，它严格介于 $f$ 在 $I$ 端点处的极限之间，
\omterm{def:g10:algebra:equation}{方程} $f(x) = k$ 在 $I$ 中有\emph{唯一}解。
\end{theorem}

\begin{proof}
存在性：取 $a, b \in I$，使 $k$ 介于 $f(a)$ 与 $f(b)$ 之间（这是可以做到的，
因为 $f$ 的取值趋近于端点处的极限），并
在 $\intcc{a}{b}$ 上应用介值定理。唯一性：若 $x < y$，由
严格\omterm{def:g12:seq:monotonic}{单调}性得 $f(x) \neq f(y)$，因此两个不同点不可能同时为
解。
\end{proof}

\begin{method}[用双射定理求解 $f(x) = k$]\label{met:g12:limcont:bijection}
要证明 $f(x) = k$ 在某\omterm{def:g10:numbers:interval}{区间}内恰有一个解：
\begin{enumerate}
  \item 作 $f$ 的\omterm{def:g10:functions:table}{变化表}（$f'$ 的符号，见
        \cref{ch:g12:deriv}）；
  \item 在 $f$ \omterm{def:g12:limcont:continuity}{连续}且严格\omterm{def:g12:seq:monotonic}{单调}的每个\omterm{def:g10:numbers:interval}{区间}上，
        比较 $k$ 与端点处的\omterm{def:g11:func:function}{函数}值或极限；
  \item 在每个这样的\omterm{def:g10:numbers:interval}{区间}上应用双射定理，并计数
        解的个数。
\end{enumerate}
借助计算器或\omterm{thm:g12:limcont:ivt}{二分法}，即可将每个解精确定位到所需的任意精度。
\end{method}

\begin{example}\label{ex:g12:limcont:cubic}
\omterm{def:g10:algebra:equation}{方程} $x^3 + x = 1$ 有唯一\omterm{def:g10:numbers:sets}{实数}解。事实上
$f(x) = x^3 + x$ 在 $\R$ 上\omterm{def:g12:limcont:continuity}{连续}且严格\omterm{def:g12:seq:monotonic}{递增}（作为
严格\omterm{def:g11:func:monotone}{递增函数}的和），$\lim_{-\infty} f = -\infty$ 且
$\lim_{+\infty} f = +\infty$，因此可用双射定理，取 $k = 1$。
由于 $f(0.6) = 0.816 < 1$ 且 $f(0.7) = 1.043 > 1$，解位于\omterm{def:g10:numbers:interval}{区间}
$\intoo{0.6}{0.7}$ 内。
\end{example}

\section{练习}

\begin{exercise}[$\star$]\label{exo:g12:limcont:1}
求下列极限：
\[
\lim_{x\to+\infty} \frac{2x^3 - x + 1}{5x^3 + x^2}, \qquad
\lim_{x\to-\infty} \bigl(x^2 + 3x - 1\bigr), \qquad
\lim_{x\to 2^+} \frac{x+1}{x-2}, \qquad
\lim_{x\to+\infty} \frac{\cos x}{x}.
\]
\end{exercise}

\begin{exercise}[$\star$]\label{exo:g12:limcont:2}
设 $f(x) = \dfrac{2x^2 - 3x + 1}{x - 1}$，其中 $x \neq 1$。
\begin{enumerate}
  \item 证明：$f(x) = 2x - 1$ 对所有 $x \neq 1$ 成立。
  \item 由此求 $\lim\limits_{x \to 1} f(x)$。能否将 $f$ 延拓为一个
        在 $1$ 处\omterm{def:g12:limcont:continuity}{连续}的\omterm{def:g11:func:function}{函数}？
\end{enumerate}
\end{exercise}

\begin{exercise}[$\star$]\label{exo:g12:limcont:3}
求曲线
$f(x) = \dfrac{3x - 1}{x + 2}$ 的\omterm{def:g12:limcont:asymptote}{渐近线}，并画出其草图。
\end{exercise}

\begin{exercise}[$\star\star$]\label{exo:g12:limcont:4}
计算
\[
\lim_{x\to+\infty} \bigl(\sqrt{x^2 + x} - x\bigr)
\qquad\text{和}\qquad
\lim_{x\to 0} \frac{\sqrt{1+x} - 1}{x}.
\]
\end{exercise}

\begin{exercise}[$\star\star$]\label{exo:g12:limcont:5}
设 $f(x) = x + \sin x$。
\begin{enumerate}
  \item 证明：$f(x) \to +\infty$，当 $x \to +\infty$ 时，尽管 $f$ 在任一
        \omterm{def:g10:numbers:interval}{区间} $\intco{A}{+\infty}$ 上都不\omterm{def:g12:seq:monotonic}{单调}。
  \item 证明：$g(x) = \dfrac{x + \sin x}{x - \sin x}$ 趋于 $1$，
        当自变量趋于 $+\infty$ 时。
\end{enumerate}
\end{exercise}

\begin{exercise}[$\star\star$]\label{exo:g12:limcont:6}
  证明\omterm{def:g10:algebra:equation}{方程} $x^3 - 3x + 1 = 0$ 恰有三个\omterm{def:g10:numbers:sets}{实数}解，
  并为每个解给出一个包含它、长度为 $1$ 的\omterm{def:g10:numbers:interval}{区间}。
  （提示：研究 $f(x) = x^3 - 3x + 1$ 的变化情况；其\omterm{def:g11:deriv:derivative}{导数}为
  $3x^2 - 3$。）
\end{exercise}

\begin{exercise}[$\star\star$]\label{exo:g12:limcont:7}
  设 $f \colon \intcc{0}{1} \to \intcc{0}{1}$ 为\omterm{def:g12:limcont:continuity}{连续}\omterm{def:g11:func:function}{函数}。证明 $f$ 有
  一个\emph{\omterm{pb:g10:functions:1}{不动点}}：存在 $c \in \intcc{0}{1}$，使得 $f(c) = c$。
  （提示：对 $g(x) = f(x) - x$ 应用介值定理。）
\end{exercise}

\begin{exercise}[$\star\star\star$]\label{exo:g12:limcont:8}
  一名徒步者沿一条山路登山，8:00 出发，20:00 到达。
  次日她沿同一条山路下山，同样 8:00 出发，
  20:00 到达。证明存在一天中的某一时刻，她在两天里
  在山路上处于完全相同的位置。（提示：引入两个
  位置\omterm{def:g11:func:function}{函数}之差。）
\end{exercise}

\begin{exercise}[$\star\star\star$]\label{exo:g12:limcont:9}
  设 $f(x) = x^5 + x^3 + x - 1$。
\begin{enumerate}
  \item 证明 $f$ 有唯一的实零点 $\alpha$，且
        $\alpha \in \intoo{0}{1}$。
  \item 描述一个用\omterm{thm:g12:limcont:ivt}{二分法}计算 $\alpha$ 的过程，使误差
        不超过 $10^{-6}$，并给出从
        $\intcc{0}{1}$ 出发所需的迭代次数。
\end{enumerate}
\end{exercise}

=\section{问题：不沾水就过不了河}

=\begin{problem}[{周末问题——介值定理
  保证了不动点、等温对径点与稳定的桌子：
  三个听起来不可思议的定理，源于同一个想法}]
\label{pb:g12:limcont:1}
  把法国地图揉成一团，随手扔到法国境内任一处：地图上
  总有一点\emph{恰好}位于它所代表的地点的
  正上方。任一时刻，赤道上两个对径点
  的温度\emph{恰好}相同。又有一张
  在凹凸不平地面上晃动的方桌，总可以
  通过\emph{旋转}使其稳定。三个“趣味戏法”，同出一个定理：\omterm{def:g12:limcont:continuity}{连续}量
  从负值变到正值时必然取到零
  （\cref{thm:g12:limcont:ivt}）。本题先训练极限，再构造
  上述三个奇妙结论，并清楚地知道定理究竟
  不保证什么。

\medskip
\noindent\textbf{第一部分——极限运算熟练度。}
\begin{enumerate}
  \item 计算 $\lim_{x \to +\infty} \dfrac{3x^2 - x}{x^2 + 5}$，
        \ $\lim_{x \to +\infty} \dfrac{2x + 1}{x^2 + 1}$，\ 以及
        $\lim_{x \to +\infty} (x^3 - x^2)$
        （\cref{met:g12:limcont:limits}）。
  \item 求
        $f(x) = \dfrac{3x - 1}{x + 2}$
        的\omterm{def:g12:limcont:asymptote}{渐近线}（\cref{def:g12:limcont:asymptote}），计算
        相关极限。
  \item 计算 $\lim_{x \to 3} \dfrac{x^2 - 9}{x - 3}$
        （用\omterm{def:g10:algebra:expand}{因式分解}）及
        $\lim_{x \to 0} \dfrac{\sqrt{x + 1} - 1}{x}$
        （用共轭式）。
  \item 由夹逼定理
        （\cref{thm:g12:limcont:squeeze}）：计算
        $\lim_{x \to 0} x \sin\frac1x$ 及
        $\lim_{x \to +\infty} \frac{\sin x}{x}$。
  \item 利用复合与主导项：
        $\lim_{x \to +\infty} \dfrac{\sqrt{4x^2 + x}}{x}$。
\end{enumerate}

\medskip
\noindent\textbf{第二部分——根的严格验证。}
\begin{enumerate}[resume]
  \item 证明\omterm{def:g10:algebra:equation}{方程} $x^3 + x = 5$ 在
        $\intcc{1}{2}$ 内至少有一个解。
  \item 证明该解在 $\R$ 上唯一
        （\cref{thm:g12:limcont:bijection}）。
  \item 证明经典结论：\emph{每个奇数次多项式
        至少有一个实根}。（写出论证
        对 $p(x) = x^3 + a x^2 + b x + c$，按主导项计算两个
        无穷极限，然后在
        一个大\omterm{def:g10:numbers:interval}{区间}上应用介值定理。）
  \item 用\omterm{thm:g12:limcont:ivt}{二分法}定位第 6 题的根：需要对
        $\intcc{1}{2}$ 作多少次二分，才能保证误差小于
        $10^{-6}$（见\cref{exo:g12:limcont:9}）？对比
        牛顿法的位数翻倍
        （见\cref{pb:g11:deriv:1}），大约需要
        多少步？
  \item 守护定理的反例：
        $g(x) = \frac1x$ 满足 $g(-1) = -1 < 0$ 且
        $g(1) = 1 > 0$，但从不为零。哪一条介值定理的假设
        不成立——该例对
        检验假设有何启示？
\end{enumerate}

\medskip
\noindent\textbf{第三部分——三个听起来不可思议的
定理。}
\begin{enumerate}[resume]
  \item 线段上的\omterm{pb:g10:functions:1}{不动点}定理：若
        $f : \intcc{0}{1} \to \intcc{0}{1}$ \omterm{def:g12:limcont:continuity}{连续}，证明
        存在 $c$ 使得 $f(c) = c$。（研究
        $g(x) = f(x) - x$ 在两个端点处的取值。）
  \item 将第 11 题转化为揉皱地图的情形：什么
        充当 $f$ 的角色，为何它会落在
        $\intcc{0}{1}$ 中（形象地说，即落在法国境内），以及
        恰好位于其自身位置正上方的点是什么？（一维情形
        在本阶段不加证明地承认——诚实的
        二维陈述为布劳威尔定理，
        详见大学分册。）
  \item 赤道：设 $T$ 为圆周上的一个\omterm{def:g12:limcont:continuity}{连续}温度，
        视为周期为 $2\pi$ 的\omterm{def:g11:func:function}{函数}，且满足
        $T(0) = T(2\pi)$。令 $g(t) = T(t) - T(t + \pi)$。
        计算 $g(0) + g(\pi)$，由此推出 $g$
        在 $\intcc{0}{\pi}$ 上某处为零，并叙述
        关于对径点的定理。
  \item 晃动的桌子：一张完全刚性的方桌立在
        \omterm{def:g12:limcont:continuity}{连续}（无台阶！）的凹凸地面上；三条桌腿
        始终接触地面，另一条腿悬空，其悬空高度 $h(\theta)$
        \omterm{def:g12:limcont:continuity}{连续}依赖于桌子的旋转角
        $\theta$。解释为何旋转四分之一圈必使
        悬空间隙变号，并由此推出
        存在某个角度 $\theta^*$，使四条腿都稳定着地。（此
        论证于 2005 年作为一个真正的定理发表；
        咖啡馆老板早就知道这一诀窍。）
  \item \cref{exo:g12:limcont:8}中的徒步者用的是
        同一方案。请叙述第 11、13、
        14 题与徒步者共用的模板：构造哪个\omterm{def:g11:func:function}{函数}，
        检验端点处的情况，再调用什么定理——用三行写出。
  \item 介值定理\emph{不}能给出什么：存在性是白得的，
        但它留下了哪两个问题未解决，以及
        本题中哪些工具回答了这两个问题（第 7 题与第 9 题）？
\end{enumerate}

\medskip
\noindent\textbf{第四部分——带\omterm{def:g12:limcont:asymptote}{渐近线}的曲线描绘。}
\begin{enumerate}[resume]
  \item 将 $f(x) = \dfrac{3x - 1}{x + 2}$ 改写为
        $3 - \dfrac{7}{x + 2}$，并利用它描述该曲线的全貌：
        \omterm{def:g12:limcont:asymptote}{渐近线}、两侧的变化情况，
        草图。（一条平移的\omterm{def:g11:func:reference}{双曲线}——
        \cref{pb:g10:reffunc:1}中的平均成本曲线也是一条。）
  \item \omterm{def:g11:func:function}{函数} $f(x) = \dfrac{x^2 - 1}{x - 1}$ 在
        $1$ 处无定义。给 $f(1)$ 赋一个什么值，才能使
        它在该点\omterm{def:g12:limcont:continuity}{连续}？分析学家为何把这类
        点称为\emph{可去}间断点？
  \item 取整\omterm{def:g11:func:function}{函数} $x \mapsto \lfloor x \rfloor$
        （满足 $\leq x$ 的最大\omterm{def:g10:numbers:sets}{整数}）：它在何处\omterm{def:g12:limcont:continuity}{连续}，
        在何处跳跃，以及为何介值定理的结论
        对它不成立（找出 $a < b$，使得
        $\lfloor a \rfloor < \frac12 < \lfloor b \rfloor$，但
        \omterm{def:g10:algebra:equation}{方程} $\lfloor x \rfloor = \frac12$ 无解）？
  \item 结语——\omterm{def:g12:limcont:continuity}{连续}性的两个侧面：局部承诺
        （\omterm{def:g11:func:function}{函数}值 $=$ 极限：不能瞬移），以及整体性的力量
        （介值定理：每个中间值都会被取到）。请将第三部分的
        四个奇妙结论归入正确的一面，
        并给出“不沾水就过不了河”这一口号精确的数学名称。
\end{enumerate}
\end{problem}
